\section{Week 7}
\sol{7.1}{(a)~The new graph has the five edges $ab$, $bc$, $ca$, $cd$, and $bd$. Reading off the neighbors of each vertex gives
\[
\deg(a)=2,\qquad\deg(b)=3,\qquad\deg(c)=3,\qquad\deg(d)=2,
\]
since $a$ is joined to $b,c$; $b$ to $a,c,d$; $c$ to $a,b,d$; and $d$ to $b,c$. The degree sum is $2+3+3+2=10=2\cdot5$, twice the number of edges, as Theorem~\ref{thm:handshake} requires. The vertices of odd degree are $b$ and $c$, so there are two of them, an even number, as Proposition~\ref{ch07:prop:odd} requires.

(b)~Deleting $cd$ from $G$ leaves the three edges $ab$, $bc$, and $ca$. Now
\[
\deg(a)=2,\qquad\deg(b)=2,\qquad\deg(c)=2,\qquad\deg(d)=0.
\]
The vertex $d$ is still part of the graph, but it has no neighbors, so its degree is $0$. The degree sum is $2+2+2+0=6=2\cdot3$, again twice the number of edges. No vertex has odd degree, so the number of vertices of odd degree is $0$, which is even.

In both parts, changing a single edge changed the degrees of exactly its two endpoints, each by one. So the degree sum changed by exactly $2$, and the parities of exactly two degrees changed. In part~(a), $b$ became odd and $d$ became even, so the count of odd-degree vertices stayed at two; in part~(b), $c$ and $d$ both became even, so the count dropped from two to zero.}
\sol{7.2}{Recall from Section~\ref{ch05:sec:euclid} that the remainder of an integer on division by four is one of the four numbers $0,1,2,3$; this holds for negative integers too, for example $-7=(-2)\cdot4+1$ has remainder $1$. Use these four remainders as boxes, and place each index $i\in\{1,2,3,4,5\}$ in the box of the remainder of $x_i$. Five indices are placed in four boxes, and $5>4\cdot1$, so by the pigeonhole principle (Proposition~\ref{ch07:prop:pigeonhole} with $k=1$) some box contains at least two different indices $i\ne j$. Then $x_i$ and $x_j$ leave the same remainder.

We placed the indices, not the values, in the boxes. This matters when two of the integers are equal: if, say, $x_1=x_2=7$, the indices $1$ and $2$ land in the same box, and the conclusion holds with $i=1$ and $j=2$. The argument covers this case without any separate treatment.

Four integers are not enough: $0,1,2,3$ leave the four different remainders $0,1,2,3$, so no two of them share a remainder. Here the pigeonhole principle does not apply, because $4$ is not larger than $4\cdot1$.}
\sol{7.3}{Take the directed cycle $a\to b\to c\to a$. It is an oriented graph: there are no loops, and each pair of vertices is joined in only one direction. It is nonempty and finite. Every vertex has an out-neighbor, since $a\to b$, $b\to c$, and $c\to a$, so there is no sink. The only hypothesis of Theorem~\ref{thm:sink} that fails is acyclicity, because $a,b,c,a$ is a directed cycle.

No oriented graph with fewer vertices lacks a sink. With one vertex there can be no edges, because a directed edge needs two different vertices; so the vertex is a sink. With two vertices $u$ and $v$, an oriented graph has at most one edge between them. If there is no edge, both vertices are sinks. If the edge is $u\to v$, then the only possible out-neighbor of $v$ would be $u$, and $v\to u$ is forbidden because $u\to v$ is present; so $v$ is a sink. (The case $v\to u$ is the same with the names exchanged.) Hence three vertices is the smallest possible number, and the directed cycle on three vertices is a smallest counterexample.

In the proof of Theorem~\ref{thm:sink}, a longest directed path in this graph is, for example, $a\to b\to c$; no directed path can have more than three vertices. Its last vertex $c$ has the out-neighbor $a$, which lies on the path. This is the second case of the proof: the arrow $c\to a$ closes the directed cycle $a\to b\to c\to a$. The proof reached a contradiction in that case only by using acyclicity. Without that hypothesis there is no contradiction, and the argument cannot conclude that $c$ is a sink.}
\sol{7.4}{The graph has no directed cycle, since every arrow leads from $a$ to $b$ or $c$, or from $b$ or $c$ to $d$, so no route returns to its start.

At $a$, the first out-neighborhood is $N_1^+(a)=\{b,c\}$. The two-step routes from $a$ are $a\to b\to d$ and $a\to c\to d$. Both end at $d$, which is different from $a$ and is not a first out-neighbor of $a$, so $N_2^+(a)=\{d\}$. The vertex $d$ is one member of this set, even though two routes reach it, so $|N_2^+(a)|=1$.

At $b$, we have $N_1^+(b)=\{d\}$. The vertex $d$ has no out-neighbors, so there is no two-step route from $b$, and $N_2^+(b)=\varnothing$. The vertex $c$ behaves in exactly the same way: $N_1^+(c)=\{d\}$ and $N_2^+(c)=\varnothing$. Finally, $d$ is a sink, so $N_1^+(d)=N_2^+(d)=\varnothing$.

Comparing sizes: at $a$ we get $1<2$, at $b$ and at $c$ we get $0<1$, and at $d$ we get $0\ge0$. So only $d$ satisfies $|N_2^+(v)|\ge|N_1^+(v)|$. This agrees with Proposition~\ref{ch07:prop:acyclic-seymour}: the graph is acyclic, and its sink $d$ is a vertex with the required property.}
\sol{7.5}{By the handshake theorem, the $n$ degrees add up to $2m$, so their average is $2m/n$. By the second half of the averaging principle (Proposition~\ref{ch07:prop:average}), some vertex has degree at most this average. Written out directly: if every degree were strictly greater than $2m/n$, adding these $n$ inequalities would give a degree sum strictly greater than $n\cdot(2m/n)=2m$, which contradicts the handshake theorem.

The conclusion is that \emph{some} degree is at most the average; it says nothing about how far the other degrees may be from it. For an example, let $n\ge3$ and take a \emph{star}: one vertex, the center, is joined to each of the other $n-1$ vertices, called leaves, and there are no other edges. The star has $m=n-1$ edges, so its average degree is
\[
\frac{2(n-1)}{n}=2-\frac2n,
\]
which is less than $2$. Every leaf has degree $1$, which is below the average, since $2-2/n\ge2-2/3=4/3>1$ for $n\ge3$. The center, however, has degree $n-1$. For $n=101$, the average degree is $200/101\approx1.98$, while the center has degree $100$. As $n$ grows, the center's degree grows without bound, while the average stays below $2$. So knowing the average does not prevent an individual degree from being very far from it.}
\sol{7.6}{No. Take the tournament on the vertices $a,b,c$ with $a\to b$, $a\to c$, and $b\to c$. Each of the three pairs is joined by exactly one directed edge, so this is a tournament. The vertex $a$ is a king, because it reaches both other vertices in one step.

Now compute the neighborhoods of $a$. Its first out-neighborhood is $N_1^+(a)=\{b,c\}$. The only two-step route from $a$ is $a\to b\to c$, since $c$ has no out-neighbors. That route ends at $c$, which is already a first out-neighbor of $a$, so $c$ is excluded and $N_2^+(a)=\varnothing$. Thus $|N_2^+(a)|=0<2=|N_1^+(a)|$, and the king $a$ fails the inequality.

The two properties measure different things. Being a king means that every other vertex is within directed distance two. The inequality compares how many vertices are at directed distance exactly one with how many are at directed distance exactly two. A king that beats almost everyone directly has a large first out-neighborhood and leaves few vertices for the second. Some vertex of this tournament does satisfy the inequality, as Proposition~\ref{ch07:prop:acyclic-seymour} guarantees, because the tournament is acyclic: its sink $c$ satisfies $0\ge0$.}
